\chapter{光源与光的传播}\label{ch:g7:light-sources-propagation}

你为什么能读这一页？灯会发光——可这一页
并非灯，但你的眼睛从每一行印刷字都能接收到光。
在能产生光的东西与只是转发光的东西之间，
存在一种光学离不开的区分；
借助这一区分，就能完整回答本书中
最古老的问题：要看见一个东西，需要什么？

\section{发光者与转发者}

\begin{definition}[自身发光的光源]\label{def:g7:light-sources-propagation:primary}
\emph{自身发光的光源}\index{自身发光的光源}自行发光，
例如太阳与恒星、火焰、炽热的灯丝、
炉中炽热的钢铁、萤火虫、闪电。它的光
不依赖其他任何发光的物体而存在——即便关掉
世间所有的灯，它仍会发光。
\end{definition}

\begin{definition}[漫射物体]\label{def:g7:light-sources-propagation:diffusion}
其他一切可见的东西都是\emph{漫射
物体}\index{漫射物体}：它接收光，再把光发散出去——
\emph{漫射}——把其中一部分射向各个方向。纸页、
苹果、月亮、你自己的脸：它们都不会自己发光；
每一个都会把接收到的光向四面八方慷慨地散射。
在黑暗中，漫射物体根本看不见。
\end{definition}

\begin{omfigure}
\begin{tikzpicture}[scale=0.8]
  % the lamp, a primary source
  \fill[omThm] (0.7,2.7) circle (0.3);
  \foreach \a in {0,45,...,315} {
    \draw[omThm, thick] (0.7,2.7) ++(\a:0.42) -- ++(\a:0.2);
  }
  \node[above, font=\small] at (1.1,3.35) {灯：自身发光的光源};
  % its light travels to the apple
  \draw[-Stealth, omThm, thick] (1.35,2.75) -- (3.55,2.15);
  \draw[-Stealth, omThm, thick] (1.35,2.45) -- (3.45,1.65);
  \draw[-Stealth, omThm, thick] (1.25,2.15) -- (3.35,1.15);
  % the apple scatters it to every side
  \draw[thick, fill=omMeth, fill opacity=0.5] (5.2,1.2) circle (0.5);
  \foreach \a in {15,60,105,150,195,240,285,330} {
    \draw[-Stealth, omDef, thick] (5.2,1.2) ++({\a}:0.6)
      -- ++({\a}:1.05);
  }
  \node[below, font=\small] at (5.2,-0.5) {苹果：漫射物体};
  \node[right, font=\small, omDef] at (7.35,2.15)
    {漫射光射向各个方向：};
  \node[right, font=\small, omDef] at (7.35,1.55)
    {房间里每只眼睛都能得到一些};
\end{tikzpicture}

{\small 分工：灯发出光；苹果
接收到光后，把一部分光射向各个方向——这就是为什么
桌旁的每个人看到的都是同一个苹果。}
\end{omfigure}

\begin{example}[漫射不是镜面反射]\label{ex:g7:light-sources-propagation:notmirror}
镜子与纸页都会把接收到的光送回去——但方式相反。
镜子沿\emph{一个}有序的方向反射，
方向（等角相等——你以前学过的反射规律），保留了
图像；纸页的微观粗糙使光
\emph{向各处}散射，破坏了图像，却让每只眼睛都能同时
接收到它。这就是为什么镜子能映出灯，而
纸页只会显示它自己，无论你站在哪里。
\end{example}

\begin{example}[月球的案子已结]\label{ex:g7:light-sources-propagation:moon}
童年时的结论如今有了正式的说法：月球
不是自身发光的光源，而是行星大小的漫射物体——
灰色的岩石把太阳光射向各个方向，
包括射向我们的方向。因此月光是二手的阳光，所谓“暗
面”，只是此刻没有光可转发的那一半，
即此刻没有光可再发出的那一半。
\end{example}

\section{看见的条件}

\begin{proposition}[视觉的两个条件]\label{prop:g7:light-sources-propagation:vision}
眼睛能看见一个物体，当且仅当
\begin{enumerate}
  \item 光\emph{离开}物体——若物体是
        自身发光的光源，则是它自身发出的光；若是转发者，则是漫射出的光——并且
  \item 这束光沿一条笔直、无遮挡的直线
        射入眼睛。
\end{enumerate}
任一条件不满足——漫射物体未被照亮，或者视线上
横着一堵墙——物体就看不见。本书中的每一个
可见性问题，都可以拆解为这两项检验。
\end{proposition}

\begin{example}[被看见的光束]\label{ex:g7:light-sources-propagation:beam}
在洁净的空气中，穿过房间的手电筒光束从侧面看不见——
空气几乎不转发任何光，光束中的光不会
转向你的眼睛。抖动一块布满灰尘的布：光束
立刻显现出来。每一粒尘埃都是一个微小的漫射体，接住光束中的光，
并把其中一部分散射向你——成千上万的尘点
为你勾勒出光束的直线。雾、薄雾和烟
为灯塔光束和电影放映机的光束起着同样的作用。
\end{example}

\section{物质与光的三种关系}

\begin{definition}[透明、半透明与不透明]\label{def:g7:light-sources-propagation:coats}
物质与光相遇有三种方式：
\begin{enumerate}
  \item \emph{透明的}\index{透明的}：光穿过
        几乎不受干扰，保持原有次序——如清澈的玻璃、
        静水、空气；图像可完整通过；
  \item \emph{半透明的}\index{半透明的}：光穿过，
        但光线已被打乱——如毛玻璃、描图纸、
        薄窗帘；亮度能通过，图像则不能；
  \item \emph{不透明的}\index{不透明的}：光被阻挡——
        被吸收或被反向漫射——如木头、石头、金属、
        你的手；不透明物体的后方是影子。
\end{enumerate}
\end{definition}

\begin{method}[用手电光给材料分类]\label{met:g7:light-sources-propagation:sort}
暗室、手电筒、一页印刷品，以及两者之间的样品：
\begin{enumerate}
  \item 把样品放在印刷品上方，手电筒位于样品
        后方；
  \item 透过样品能读出字迹吗？ --- 透明的；
  \item 能看到亮光却看不清字吗？ --- 半透明的；
  \item 看到一片黑暗，且后方有影子吗？ --- 不透明的。
\end{enumerate}
认真检验：很多材料会随厚度改变类别 --- 薄薄一层
水是透明的，深海却漆黑如夜；一张
纸是半透明的，合上的书则不透明。
\end{method}

\begin{omfigure}
\begin{tikzpicture}[scale=0.8]
  \foreach \x/\name in {0.5/透明, 4.4/半透明,
      8.3/不透明} {
    \fill[omThm] (\x,1.5) circle (0.22);
    \draw[thick] (\x+1.3,0.4) rectangle (\x+1.7,2.6);
    \node[below, font=\small] at (\x+1.6,-0.0) {\name};
  }
  \draw[-Stealth, omThm, thick] (0.75,1.5) -- (1.78,1.5);
  \draw[-Stealth, omThm, thick] (2.22,1.5) -- (3.4,1.5);
  \draw[-Stealth, omThm, thick] (4.65,1.5) -- (5.68,1.5);
  \foreach \a in {28,10,-12,-30} {
    \draw[-Stealth, omThm, thin] (6.12,1.5) -- ++({\a}:1.1);
  }
  \draw[-Stealth, omThm, thick] (8.55,1.5) -- (9.58,1.5);
  \draw[line width=3pt, omExo] (10.02,0.55) -- (10.02,2.45);
  \node[right, font=\footnotesize, omExo] at (10.1,0.8) {影子一侧};
\end{tikzpicture}

{\small 光线的三种去向：有序穿过；穿过但已被打乱；
被阻挡，后方有影子。}
\end{omfigure}

\begin{example}[建筑中的三类材料]\label{ex:g7:light-sources-propagation:architecture}
建造者有意使用这三类材料。透明的窗玻璃：
既能看见外景，又能透进光线。半透明的浴室窗玻璃与磨砂
办公隔断：让日光透入、保护隐私 --- 只透亮度、
不透影像，正是它们的本职。不透明的墙壁
与百叶窗：按需要营造黑暗。灯罩则是
半透明性服务于舒适的典型：炽热灯丝的刺眼眩光
被打散，变成柔和、看不见光源的光晕。
\end{example}

\begin{remark}[无光散射，便无所见]\label{rem:g7:light-sources-propagation:space}
谜语：一位宇航员在群星之间，\emph{向侧面}看去，
穿过充满阳光的虚空。她看到什么？一片漆黑 ---
尽管阳光从她眼前涌过。太空中几乎没有
值得一提的尘埃：没有任何物质把掠过的光
向她漫射，未进入眼睛的光什么也不会成像。
白天的天空之所以呈蓝色，正是因为空气\emph{确实}
将阳光向下传到我们这里 --- 整个大气层的
柔和漫射。没有空气，即使正午天空也是黑的：月球上的
宇航员站在阳光下，头顶是无星的黑色穹顶。
\end{remark}

\section{练习}

\begin{exercise}[$\star$]\label{exo:g7:light-sources-propagation:1}
将下列物体分类：烛火；\ 月球；\ 萤火虫；\ 这页纸；\ 一
个红色交通信号灯；\ 黄昏时的行星金星。
\end{exercise}

\begin{exercise}[$\star$]\label{exo:g7:light-sources-propagation:2}
漫射物体对所接收的光做了什么 ---
这些光在暗室里又会怎样？
\end{exercise}

\begin{exercise}[$\star$]\label{exo:g7:light-sources-propagation:3}
说出视觉的两个条件，并运用它们解释：为什么你
看不见一面完全干净的镜子的\emph{自身}表面，只能看到
它所反射的东西？（注意——这是个细致的问题；请简要回答：
它的光洁表面拒绝对光做什么？）
\end{exercise}

\begin{exercise}[$\star$]\label{exo:g7:light-sources-propagation:4}
为什么围坐在桌旁的每个人看到的是同一个苹果，
而只有某个幸运的眼睛位置，才能看到灯在
勺子上的反射？
\end{exercise}

\begin{exercise}[$\star$]\label{exo:g7:light-sources-propagation:5}
为什么手电筒的光束在洁净的空气中从侧面看不见，
而在薄雾中却看得见？说出这一过程的名称。
\end{exercise}

\begin{exercise}[$\star$]\label{exo:g7:light-sources-propagation:6}
判断下列物体各属于哪一类：窗玻璃；\ 描图纸；\ 砖块；\ 浅层
清水；\ 磨砂灯泡的外壳；\ 铝
箔。
\end{exercise}

\begin{exercise}[$\star$]\label{exo:g7:light-sources-propagation:7}
半透明性在浴室窗户与
灯罩中起什么作用？什么光能通过，什么被阻挡？
\end{exercise}

\begin{exercise}[$\star$]\label{exo:g7:light-sources-propagation:8}
月球被阳光照亮 --- 然而月球上的宇航员在
正午却看到黑色天空。请借助
\cref{rem:g7:light-sources-propagation:space}，并说明
为什么地球正午的天空反而是明亮的蓝色。
\end{exercise}

\begin{exercise}[$\star\star$]\label{exo:g7:light-sources-propagation:9}
一张纸在手电筒前透出柔和的暖光；
合上的书则完全挡住光线。厚度对材料的类别归属有
什么影响？请再举一个
能改变类别的例子。
\end{exercise}

\begin{exercise}[$\star\star$]\label{exo:g7:light-sources-propagation:10}
电影的小把戏：投影仪的光束在洁净的大厅里是看不见的，
但电影里却常见光束在空气中戏剧性地划过。剧组须往空气中
加入什么？他们借助的又是
哪种物理原理？
\end{exercise}

\begin{exercise}[$\star\star$]\label{exo:g7:light-sources-propagation:11}
你在篝火旁看到朋友的脸。请沿着光的完整路径——光源、漫射、
直线传播——追踪一遍，并数一数，
若你\emph{也}在湖面上看到这张脸的倒影，
那么这条路径上共有几个漫射物体？
\end{exercise}

\begin{exercise}[$\star\star\star$]\label{exo:g7:light-sources-propagation:12}
大自然自己的光影秀：晴朗的白天，远离太阳的天空呈蓝色，
太阳圆盘却白得耀眼——而日落时，低垂的太阳变成橙色，
上方的云彩则泛起粉红。仅利用“空气会使穿过它的光发生少量漫射，
且路径越长漫射越多”这一点，
解释低垂的太阳为何损失了一部分光，
以及被夺走的那部分光去了哪里。
（完整的颜色故事留给明年的光谱一章；
此处只用“路径越长，漫射越多”来推理。）
\end{exercise}

\section{问题：为博物馆布光}

\begin{problem}[{周末问题 --- 新雕塑展厅；灯光设计师的夜班；
光源、漫射物体与三种
玻璃}]\label{pb:g7:light-sources-propagation:1}
博物馆新展厅周一开放：一尊大理石雕像、墙上的画作、
一个玻璃展柜，以及一位紧张的馆长。你
协助灯光设计师工作。

\medskip\noindent\textbf{第一部分 --- 现场的基本原理。}
\begin{enumerate}
  \item 馆长问，为何雕像“白得如此耀眼”，
        仍然需要灯光照明？请
        用视觉的两个条件来回答。
  \item 灯光亮起之后，列出大厅中的
        自身发光的光源与漫射物体：灯、雕像、画作、
        参观者，以及展柜的玻璃。
  \item 大理石似乎从各个方向都在发光，
        没有刺眼的亮点。这属于哪种反射方式？
        为什么每个参观位置都能接收到它？
  \item 相比之下，一块抛光的黄铜牌匾只产生
        一处刺眼的灯光反射。请说出另一种反射方式，
        并说明这束刺眼反光射向何处所遵循的规律。
\end{enumerate}

\medskip\noindent\textbf{第二部分 --- 控制眩光。}
\begin{enumerate}[resume]
  \item 裸露的灯泡会在参观者眼中造成眩光。设计师为
        每只灯泡加上一个磨砂灯罩。磨砂灯罩属于哪一类材料，
        它对灯丝发出的光有什么影响？
  \item 画作必须打光，但从主要观赏位置看去，
        画作漆面上不得出现任何灯具的\emph{反射}像。
        利用镜面反射的规律，设计师应把灯具
        置于相对画作与观众的何处？
  \item 展柜的透明玻璃几乎看不见——
        直到参观者白色的袖子像幽灵一样映在其中。
        为什么透明玻璃既能透光，又会微微地镜面反射？
        设计师在展柜后方设置的黑色墙龛，
        起到什么作用？
  \item 灰尘微粒有时会飘进聚光灯的
        光束中，在空中勾出明亮的细线。馆长说：
        “美极了！留着吧。”设计师说：“这说明空气
        很脏。”谁对物理的判断是正确的，
        这些细线又是什么？
\end{enumerate}

\medskip\noindent\textbf{第三部分 --- 夜间测试。}
\begin{enumerate}[resume]
  \item 熄灯并封好百叶窗后，保安报告说大厅
        “漆黑一片——连白色雕像都看不见”。
        请用视觉的两个条件解释，为何白色
        在黑暗中毫无作用。
  \item 有人提议在屋顶装天窗，用磨砂玻璃覆盖。
        白天它能提供什么，又不能提供什么
        （如云彩的景象、直射阳光造成的眩光）？
  \item 馆长担心阳光会使画作褪色，
        因此要求“不含任何阳光的光”。设计师笑着说：
        白天，即使是朝北窗户透入的柔和光，\emph{其实}也是阳光。
        请解释：是什么把它“转送”过来的？
  \item 写出设计师的报告：三句话——一句谈为何漫射是
        博物馆的好帮手，一句谈镜面反射须留意的位置，
        一句谈为每种用途所选玻璃的
        类别。
\end{enumerate}
\end{problem}
